\subsection{Summary of contributions}
We extended a computational sCPG model --- a two-legged half-centre
network in the Hodgkin--Huxley lineage
running from Rybak 2006 through Shevtsova 2015 and Danner 2017 to
Zhang 2022~\cite{rybak2006,shevtsova2015,danner2017,zhang2022} --- with four
mechanisms absent from that lineage: (i)
STDP that lets the network tune its own
sensory-to-rhythm-generator synapses rather than have them hand-set; (ii) a
closed loop from muscle-force feedback (Ia) into the reciprocal-inhibition
core; (iii) motoneuron pools driving explicit muscle proxies with force and
length dynamics; and (iv) an externally paced cutaneous drive that emulates
heel-to-toe pressure transfer during stance. Using this model we asked how a
spinal CPG's own synaptic weights should be organised to produce a stable
walking rhythm, and how that self-tuned rhythm degrades as sensory feedback
is progressively withdrawn --- the central manipulation in rodent and human
studies of SCI rehabilitation.

Three findings anchor the paper. First, the three plastic synapses --- the
cutaneous CUT$\to$RG-E projection and the homonymous proprioceptive
projections Ia-E$\to$RG-E, Ia-F$\to$RG-F --- converge to the same steady
strength regardless of where they start, of which leg they are in, and of
walking speed (\autoref{sec:results:weights}); the network organises its own
descending gain rather than requiring it to be prescribed by the modeller.
Second, across the full grid of five locomotion modes and five learning
rates the quality of the alternating gait is not maximised by the fastest
learning, but peaks at an intermediate rate and falls off at both extremes
(\autoref{sec:results:comparison}). Third, removing load-dependent sensory
feedback and removing the cutaneous pacing drive have dissociable effects:
a model with only the pacing drive intact keeps stepping right through full
unloading, while a model in which both proprioceptive and cutaneous
feedback are gated by load collapses --- reproducing, in a spiking
closed-loop network, the clinical dissociation between unassisted and
epidurally-stimulated stepping after motor-complete SCI
(\autoref{sec:results:epidural}).

\subsection{Relation to prior computational and experimental sCPG work}
Zhang et al.\ (2022)~\cite{zhang2022} name two limitations of their own
four-rhythm-generator computational sCPG model explicitly: it has no
biomechanics or sensory feedback, and it has no motoneurons, treating
rhythm-generator output as a direct stand-in for motor output. Every
synaptic weight in their model, like its Rybak/Shevtsova/Danner
predecessors, is hand-tuned. Our model addresses both limitations for a
single hindlimb pair: the Ia force/length feedback loop supplies the
missing sensory closed loop, and explicit motoneuron pools driving a
saturating muscle model supply the missing motor periphery. On top of this
we add a mechanism none of these models have --- activity-dependent
plasticity on the afferent synapses --- so that the descending and
cutaneous gains are discovered by the network rather than fixed by the
experimenter.

The trade-off is architectural detail. Zhang et al.\ resolve genetically
identified commissural and long-propriospinal interneuron classes (V0V,
V0D, V1, V2a, V3, Shox2) across four limbs and reproduce V3-silencing
experiments quantitatively; our model collapses the commissural pathway to
two hand-set inhibitory projections between homologous half-centres and
covers only one limb pair. We also do not reproduce gait-transition
phenomena (walk $\to$ trot $\to$ gallop $\to$ bound): our gait is
externally paced rather than emergent from a single tonic-drive parameter.
The two efforts are complementary rather than competing: theirs
demonstrates architectural realism at the cost of a closed loop; ours
demonstrates a closed, self-organising loop at the cost of architectural
detail. A natural synthesis --- hand-curated interlimb architecture wrapped
in a plastic, closed sensorimotor loop --- is a direct extension of both.

Our loading manipulation and its epidural-stimulation analogue are drawn
directly from a second, experimental lineage: the rat and human
spinal-cord-injury work of Courtine, Lavrov, Edgerton, Gerasimenko and
colleagues showing that epidural stimulation of the lumbosacral spinal
cord can recover stepping across a graded series of load-bearing
conditions, from fully weight-supported walking to hindlimb
suspension~\cite{lavrov2008,edgerton2008,courtine2009,harkema2011}. That
literature establishes the phenomenon --- rhythmic electrical stimulation
of the cord restores stepping when descending drive is lost, in a
load-dependent manner --- largely from kinematic and electromyographic
readouts. Our model does not simulate the stimulation itself, but it
offers a mechanistic, spiking-network account of one candidate route to
the same phenomenon: a closed sensorimotor loop in which a rhythmic
afferent volley, held constant across loading, substitutes for the
load-gated sensory drive that unassisted stepping loses
(\autoref{sec:results:epidural}). In that sense the model is positioned to
connect the computational sCPG lineage above with this experimental
SCI-rehabilitation lineage, and we return to that connection in the
Clinical implications below.

\subsection{A shared mechanism underlies the learning-rate and loading
optima}
Two of our findings share the same underlying mechanism, which is worth
making explicit before turning to its clinical reading. Learning only
improves the gait once it has had time to organise the cutaneous weight
appropriately, and only within a specific range of final weight: at the
fastest rate ($\lambda=10^{-2}$) the descending weight overshoots useful
values and the alternation becomes slightly noisier, while at the slower
rates ($\lambda=10^{-5}, 10^{-6}$) the weight never gets there at all, so
the comparison across learning rates (\autoref{sec:results:comparison})
shows a broad optimum around $\lambda=10^{-3}$--$10^{-4}$ rather than a
monotonic benefit of faster learning. The same logic explains the
non-monotonic recovery we observe from toe to air stepping at slow learning
rates: an intermediate, partially-developed weight is poorly enough timed
to \emph{interfere} with the paced core's own alternation, whereas a fully
withdrawn afferent drive (air stepping) simply leaves the paced core to run
on its own, unperturbed floor. In both the learning-rate axis and the
loading axis, then, \emph{some} plastic drive organised correctly is better
than none, but a large, poorly-organised drive can be worse than either
extreme.

\subsection{Clinical implications}
The present computational approach yields several insights for the
clinical rehabilitation of individuals with SCI. Apart from recapitulating
locomotor behaviour, the proposed computational framework provides a
mechanistic account of how sensorimotor interactions and activity-dependent
plasticity might enable functional recovery. First, sensory inputs are
critical for facilitating locomotor rehabilitation, not as a replacement
for lost supraspinal descending input, but as a catalyst for plastic
adaptation in the spinal cord sensorimotor networks.

Our simulation results show that, when appropriately controlled through
load and/or activity, such sensorimotor input drives spinal adaptation and
subsequently restores coordinated locomotor outputs
(\autoref{sec:results:weights}).

Weight-bearing conditions fostered increasing potentiation of cutaneous
input and stabilisation of locomotor output, whereas reduced support
reduced sensory amplification and led to a more disorganised gait
(\autoref{sec:results:comparison}). These findings parallel clinical
observations that repetitive task-specific locomotor rehabilitation ---
body-weight-support treadmill walking, robot- and manually-assisted
stepping --- restores stepping in individuals with motor-incomplete SCI by
engaging residual plasticity in the spinal sensorimotor circuitry rather
than by regenerating damaged descending
pathways~\cite{wernigMuller1992,dietz1998,behrmanHarkema2000}. The model
provides a biologically plausible hypothesis for the benefits of epidural
stimulation in SCI. Reduced proprioceptive and exteroceptive inputs from
unloaded conditions resulted in poor locomotor performance.

The rhythmic input was, however, continuously driven by externally
delivered rhythmic cutaneous input, and these effects resulted in restored
alternating activity between flexors and extensors
(\autoref{sec:results:epidural}), a phenomenon observed in experiments in
animal models as well as in patients undergoing epidural stimulation of
the spinal cord~\cite{lavrov2008,edgerton2008,courtine2009,harkema2011}.

\paragraph{Which injury this model corresponds to.}
The two experimental manipulations map onto different points on the
clinical severity spectrum. The plasticity-driven weight-bearing arm
(\autoref{sec:results:weights}--\autoref{sec:results:comparison}) leaves
the brainstem-to-rhythm-generator (BS$\to$RG) drive tonic and non-zero
throughout, so it structurally represents a state in which some
supraspinal descending input survives the lesion --- clinically,
motor-incomplete SCI (ASIA Impairment Scale grades C/D), the population in
which BWSTT and robot- or manually-assisted training reliably restore
stepping~\cite{wernigMuller1992,dietz1998,behrmanHarkema2000}. Notably, the
synapses that are plastic in this arm, CUT$\to$RG-E and Ia$\to$RG, are not
the projections a thoracic or cervical lesion actually damages: those
afferent-to-interneuron connections are intraspinal and caudal to the
lesion, and do not cross it. What STDP models here is therefore not repair
of a severed pathway, but activity-dependent strengthening of an intact,
underused sensorimotor loop that compensates for reduced descending drive
--- the mechanism BWSTT and related therapies are thought to exploit. The
unloading/ESCS arm (\autoref{sec:results:epidural}), by contrast, models a
more severe injury in which even that residual drive is functionally
inadequate and an externally imposed rhythmic afferent volley substitutes
for it outright; this is closer to the regime in which epidural
stimulation is used clinically, often in motor-complete or near-complete
SCI where descending drive is minimal or absent~\cite{harkema2011}.

From this perspective, ESCS could function not just by increasing spinal
excitability, but by restoring or amplifying sensory signals necessary for
engaging the central pattern generators of the spinal cord to control
locomotion. Therefore, neuromodulation in SCI might be more efficient if it
is part of an integrated approach which includes not only electrical
stimulation, but also functionally relevant sensory stimuli. An additional
clinically relevant observation from our work is that the effectiveness of
recovery depends on the degree of plasticity.

Simulations reveal that intermediate levels of plasticity can produce more
robust, functional gait than very low or very high values ---
a range in which functional rehabilitation appears optimal, not maximised
(\autoref{sec:results:comparison}). The optimal levels of learning rate
would thus correspond to those that enable spinal cord motor networks to
adapt while avoiding extreme rewiring, and provide a useful heuristic guide
for determining appropriate intensities for rehabilitation and for
managing the effects of pharmacological or electrical neuromodulation
therapies on spinal adaptation. The optimal range of plasticity should be
dynamically adapted to each patient's neurophysiological state and
clinical characteristics.

\paragraph{Interpreting the learning-rate axis as plasticity capacity.}
One natural reading of the $\lambda$ axis is developmental. The capacity
for activity-dependent synaptic change is not fixed across the lifespan:
it is greatest during early postnatal critical periods and declines with
maturity, and correspondingly, neonatal and juvenile rodents recover
locomotor function after spinal injury considerably better than adults.
It is therefore tempting to read our fastest learning rates as
``younger'' and our slowest as ``older'' animals.\todo{cite developmental
critical-period / age-related decline in spinal plasticity and
age-dependent locomotor recovery after SCI} We suggest this analogy
should be treated as a qualitative ordering rather than a calibration:
$\lambda$ is a phenomenological STDP rate constant, it was not fitted to
any measurement of plasticity in animals of a given age, and age is only
one of several determinants of plasticity capacity --- time since injury,
prior activity history, and pharmacological or electrical neuromodulation
all modulate the same underlying quantity. Read this way, the $\lambda$
axis is better described as spanning \emph{plasticity capacity}, of which
age is one contributor.

With that caveat, the interpretation yields a concrete and testable
prediction. Because gait quality in our sweep peaks at intermediate
$\lambda$ rather than increasing monotonically
(\autoref{sec:results:comparison}), the model predicts that maximal
plasticity capacity is not automatically optimal for functional outcome
--- that a system rewiring too readily can be as poorly coordinated as one
that barely rewires at all. This resonates with the clinical picture in
which excessive or misdirected plasticity after SCI manifests as
maladaptive outcomes such as spasticity, aberrant sprouting and
neuropathic pain rather than improved locomotion.\todo{cite maladaptive
plasticity after SCI} The prediction is directly testable by running the
same rehabilitation protocol across age groups and asking whether
intermediate-age animals show the best coordination recovery, which would
be a stronger validation of the learning-rate axis than any of the
comparisons we are currently able to make.

Furthermore, the present computational framework might serve as a useful
tool for designing, optimising, and validating interventions before their
clinical application.

It is feasible to investigate, via simulation studies, various clinical
rehabilitation interventions (e.g., robot-assisted walking at various
assistance levels and weight-supported environments, ESCS at different
locations and frequencies, combination therapies) to identify the most
effective, non-invasive treatments. Finally, integrating this framework
with individual patient-specific data might pave the way for the
development of virtual-twin approaches that could predict, simulate, and
optimise patient outcomes after specific rehabilitation therapies in a
highly personalised fashion.

\subsection{Limitations}
In addition to the positive implications outlined above, it is important
to note some limitations of the proposed model. To simplify the biological
complexity of the spinal locomotor system, a level of abstraction was used
in defining sensorimotor pathways, interneurons and motor outputs. The
model does not include an adaptive supraspinal control mechanism (the
brainstem drive is a tonic input), detailed muscle-tendon
biomechanics beyond the saturating force/activation model used here,
fatigue, or other aspects of longer-term structural plasticity.

It also omits features of SCI such as chronic damage to the dorsal and
ventral roots, glial scarring, or the diversity of human SCI.

It should also be highlighted that the model was not calibrated on any
experimental data collected from actual individuals with SCI. Hence the
values of the STDP parameters can only be viewed as computational
representations of plasticity rather than direct measures of spinal
learning capacity. Consequently, the present model should be interpreted
as hypothesis-generating rather than predictive in nature, and further
studies incorporating more detailed muscle dynamics, more realistic
sensory transduction, supraspinal inputs, and patient-specific data will
be required to establish greater biological realism.

These conceptual limitations are compounded by several specific modelling
choices worth stating explicitly:

\begin{enumerate}
  \item \textbf{Interneuron and interlimb detail.} We resolve two hand-set
        commissural projections rather than the genetically identified
        V0/V1/V2a/V3 interneuron classes, and we model one hindlimb pair
        rather than fore--hindlimb coordination~\cite{zhang2022}; we
        likewise do not reproduce gait-transition phenomena
        (walk $\to$ trot $\to$ gallop $\to$ bound).
  \item \textbf{No stochastic replication.} Results are deterministic,
        fixed-seed point estimates; the ten-point initial-weight sweep
        (\autoref{sec:results:weights}) substitutes for repeated stochastic
        trials in establishing that a result is not an artefact of a
        particular initialisation.
\end{enumerate}

\subsection{Concluding remark}
A computational sCPG half-centre oscillator, closed in the loop with
motoneurons, muscle proxies, and load-dependent sensory feedback, tunes its
own descending gain through spike-timing-dependent plasticity to a stable
locomotor attractor --- one that is largely independent of how the network
starts, robust across the rat locomotor speed range, and best expressed at
an intermediate degree of plasticity rather than the maximum
achievable. Withdrawing load-dependent sensory feedback degrades that gait,
and restoring a rhythmic, load-independent afferent volley --- the model's
analogue of epidural stimulation --- rescues it, reproducing a central
dissociation in the SCI rehabilitation literature within a single
mechanistic, closed-loop model. Beyond the specific findings, the model
demonstrates that a spinal CPG's descending and sensory gains need not be
prescribed by the modeller: given a plastic synapse, appropriately timed
sensory drive, and enough time, the network finds them on its own. That
property, together with the model's explicit and inspectable parameters,
is what makes it a plausible substrate for exploring --- and eventually,
with patient-specific calibration, for helping to design --- sensorimotor
rehabilitation strategies after spinal cord injury.
